\documentclass{article}
\usepackage{graphicx} % Required for inserting images

\usepackage{amsmath, amsthm, amssymb, amsfonts}

\title{EDP Mutation}
\author{Marie Doumic, Guillaume Garnier}
\date{\today}
\newcommand{\p}{\partial}
\newcommand{\diff}{\,\mathrm{d}}
\newcommand{\RR}{\mathbb{R}}
\newcommand{\f}[2]{\frac{#1}{#2}}
\begin{document}

\maketitle

\section{Introduction}

Article~\cite{coville2007non}: $u$= "gene fraction of the mutant" . From Fisher
\[u_t-\Delta u=u(1-u)
\]
derives the nonlocal version with convolution kernel $J$:
\[u_t- (J*u-u)=f(u)\]
assumes symetric $J$: $J(x)=J(-x)$.  Be careful: their variable $x$
is a spatial variable and not at all a trait which is mutating ! 

For us different: we have for instance only deleterious mutations, so that $Supp(J)\subset (-\infty, 0),$ and no reason for $J$ symetric. However the symetry of $J$ is maybe not important in their proof? For us $f=0$ whereas for them they take $Supp(f)=(0,1)$ and $f>0$ on $(0,1):$ monostable case. If there is competition it should be nonlocal, of type $f(u)=u(1-\int w(x)u(x)dx)$ with $w$ a weight function - the simplest being $w(x)=1$. They look for traveling wave solutions of type $u(x,t)=U(x+ct)$ which, in the monostable case, progressively invade all the space.
They link their model to the  mutation-selection model of~\cite{fourniermeleard}, which they write as being
\[\p_t u=J*u-u +(1-K*u)u \]
and this is much more coherent for us: the competition should occur with the whole population.

Possible questions for our question: does there exist travelling wave solutions ? Does it have a meaning if $f=0$ ? In such a case did someone answer to the question: How 0 is reached ?

In~\cite{cloez2022fast} the authors study a "particular case of Kimura's replicator-mutator model" called "house of cards model of mutations": the mutation does not depend on the original trait, hence $J*u$ is replaced by $Q(x) \int u(t,y)dy.$ For us on the contrary it depends in a "self-similar" way of the original trait so that we have $J*u$ simplifying the full model which would be $J(x,\cdot) *y$.

In~\cite{burger1996stationary} the equation (1.3) seems exactly the same as ours (slightly more general since not only a convolution and allows non-poissonian mutation rate $\lambda(x)$) and it adds it to a selection process to obtain (1.5). Big bibliography.

New articles given by Pierre Gabriel: Gil Hamel Roques Martin. To see together. Here also there is a selection term and it seems unavoidable. 

Question: is it interesting to study the model without selection ? Biological justification for it ? Has it been studied in the previous literature as an intermediate step ? 

For us it seems that the "most-adapted trait" [name to check] is $+\infty$ since the higher the growth rate the higher the fitness: $\tau=\lambda$ [$\tau$ growth rate and $\lambda$ Malthus parameter]. In reality certainly not. Especially the fastest growth= the more mutation and especially the more deleterious mutations.

mutations not poissonian: increase with growth rate - HOW ? linear dependency easy to justify biologically

+ mutation kernel as in the simplest case

Papier par Burger et Bomze: notre equation est un cas particulier de leur equation 1.5.

OBTENTION OF THE EQUATION (1.5) of Burger and Bomze: Consider the equation
\begin{align*}
    \p_t u = xu(t,x) - \lambda u(t,x)+ \lambda \int_\RR u(t,y)\kappa\left(\frac{x}{y}\right) \frac{\diff y}{y}, \qquad t \in \RR_+ \,, x \in \RR
\end{align*}
For any $(t, x) \in \RR_+ \times \RR$, we define $p(t,x) = \frac{u(t,x)}{\int_\RR u(t,y)\diff y}$. 
We have (justify the interversion of $\int$ and $\p$
\begin{align*}
    \p_t p(t, x) 
    &= \frac{\p_tu(t,x) \cdot \int_\RR u(t,y)\diff y - u(t,x) \int_\RR \p_t u(t,y)\diff y}{\left( \int_\RR u(t,y)\diff y \right)^2} \\
    &= \frac{\left[ xu - \lambda u + \lambda \int_\RR u\kappa\left(\frac{x}{y}\right) \frac{\diff y}{y}\right] \cdot \int_\RR u\diff y - u\int_\RR \left[ yu - \lambda u + \lambda \int_\RR u(t,s)\kappa\left(\frac{y}{s}\right) \frac{\diff s}{s}\right]\diff y}{\left( \int_\RR u(t,y)\diff y \right)^2} \\
    &= xp - \lambda p + \lambda \int_\RR p\kappa\left(\frac{x}{y}\right) \frac{\diff y}{y}
    -  p\int_\RR \left[ yp - \lambda p + \lambda \int_\RR p(t,s)\kappa\left(\frac{y}{s}\right) \frac{\diff s}{s}\right]\diff y
\end{align*}
Since, $\int_\RR p \diff y = 1$, 
\begin{equation}\label{eq:BurgerBomze}\begin{array}{ll}
    \p_t p(t, x) 
    &= xp - \lambda p + \lambda \int_\RR p\kappa\left(\frac{x}{y}\right) \frac{\diff y}{y}
    -  p\int_\RR yp \diff y + \lambda p - \lambda \int_\RR \int_\RR p(t,s)\kappa\left(\frac{y}{s}\right) \frac{\diff s}{s}\diff y\\
    &= xp 
    + \lambda \int_\RR p\kappa\left(\frac{x}{y}\right) \frac{\diff y}{y}
    - p\int_\RR yp \diff y 
    - \lambda p \int_\RR \int_\RR p(t,s)\kappa\left(\frac{y}{s}\right) \frac{\diff s}{s}\diff y\\
    &= xp 
    + \lambda \int_\RR p\kappa\left(\frac{x}{y}\right) \frac{\diff y}{y}
    - p\int_\RR yp \diff y 
    - \lambda p\end{array}
\end{equation}
Where the last equality is obtained because 
$$\int_\RR \int_\RR p(t,s)\kappa\left(\frac{y}{s}\right) \frac{\diff s}{s}\diff y = \int_\RR \int_\RR p(t,s)\kappa\left(\frac{y}{s}\right) \diff y \frac{\diff s}{s} = \int_\RR \int_\RR p(t,s)s\kappa\left(y\right) \diff y \frac{\diff s}{s} = 1.$$
(Fubini is ok because all quantities are positive). So we have~\eqref{eq:BurgerBomze}
is equation (1.5) of Burger and Bomze but for a variable $x=exp(z)$ so that with Burger-Bomze notations we would have $w(y)=exp(y)$.

If we do a change of variable 
\subsection{Summary of the articles}

\textbf{Eshel 1972, Evolution processes with continuity of types}: Study the long range evolutionary traits in a population with an infinie numpber of types. Goal : Asymptotic behaviour. 

At each step : A offspring may born from its parent, and mutate. 

Here, the authors assume that each mutations are equi-distributed around the parental type (Equation 2.4). It's not the case with our kernel. Futhermore, the authors assume that time is discrete. Here again, it's not the case with our model. \\

\noindent\textbf{Karlin 1988, Non-Gaussian phenotypic models of quantitative}
Impossible to find this article. \\

\noindent\textbf{Bürger 1989, Unbounded One-Parameter Semigroups}: Model of mutation and selection. The selection is against all the individuals. In the paper, $m(x)$ represents the fitness of individual with trait $x$. In our model, $m(x) = x$. I didn't see hypothesis on the mutation kernel $u$, except mutations come from a convolution operator. 


They assume that $m$ satisfies :
\begin{enumerate}
    \item (m1) : $m$ is real-valued and $m \in L^1_{loc}(M, \nu)$ (ok for us)
    \item (m2) : $\text{ess sup}_{x \in M} m(x) = +\infty$ (ok for us)
    \item (m3) : $m_{-}$ is bounded on compact set (ok for us)
    \item (m4) : $e^m$ is submultiplicative (not ok for us. In the multiplicative case, we have $m(z) = e^{z}$). Definition of submultiplicative: there exists a constant $C$ such that
    \[ 
    e^{m_+(x+y)}\leq C e^{m_+(x)}e^{m_+(y)}\]
\end{enumerate}

The problem is the assumptions on their $u(x,y),$ for us it is $\kappa(\f{x}{y})\f{1}{y},$ given on their p. 86: they want $u$ either depending only on $x$ or being a convolution kernel. This is not the case for us. If we do a change of  variable $x=\exp(z)$ then our multiplicative convolution kernel becomes a convolution kernel, but then we have $m(x)=e^z$ which is no more submultiplicative. 


p.78 : They present an other work where $m = sx$ and said that no stationary solution can be found. 

The article use an approach with semi-group that the authors have used in an other paper. They proved there exist a unique solution, that depends continuously on the initial data.  

Avant de regarder l'equation avec taux de croissance en $x$ pour la population on peut regarder l'equation en lignee ou il n'y a pas de croissance de la pop:
\begin{align*}
    \p_t v =  - \lambda v(t,x)+ \lambda \int_\RR v(t,y)\kappa\left(\frac{x}{y}\right) \frac{\diff y}{y}, \qquad t \in \RR_+ \,, x \in \RR
\end{align*}

Comment resoudre le prob pour trouver un equilibre en pop ? Selon Lydia, plus le taux d'elongation est grand et plus les mutations s'accumulent. Or elles sont majoritairement deleteres donc cela pourrait creer un equilibre ?

\section{Pistes/directions de recherche}

Cas le plus simple deja ecrit: bien comprendre tout. Au moins grandes idees.

Garder une croissance proportionnelle a $x$ car c'est ce qui correspond a Lydia = taux de croissance. Avec cette contrainte, voir a quel point cela sort du champ deja etudie - il semble que les hypotheses ne sont pas les memes. La question est alors: est-ce qu'on peut modifier le taux de mutation, autant le taux total que le noyau de mutation, en autorisant aussi des noyaux non auto-similaires, de facon a avoir convergence vers un etat stationnaire (ce qui est plus realiste)? En theorie: plus on grossit vite et plus on mute, mais on peut imaginer en revanche que les mutations deleteres soient avantagees. Donc 2e modele le + simple serait $\lambda$ proportionnel a x et noyau de mutation ad hoc.
\bibliographystyle{plain}       % (uses file "plain.bst")
\bibliography{Tout_Marie}           % expects file "biblio.bib"

\end{document}
